\chapter*{Preface}
\addcontentsline{toc}{chapter}{Preface}

These are my seminar notes from reading Kodaira's lecture notes
\emph{複素多様体と複素構造の変形 I} (1968) during my undergraduate research.
The original frequently states a result and leaves its verification to the
reader; the purpose of these notes is to write out those omitted steps in full.

\section*{Conventions}

\begin{itemize}
  \item Manifolds are assumed second countable and Hausdorff unless stated otherwise.
  \item $\C,\R,\Z,\N$ denote the complex, real, integer, and natural numbers.
  \item $\dbar$ is the Dolbeault operator; $\Oh$ is the structure sheaf.
  \item $\coh{q}{X}{\mathcal{F}}$ denotes sheaf cohomology unless noted otherwise.
\end{itemize}
